% =====================================================================
%  H. Brøchner
%  Et Svar til Professor R. Nielsen  /  A Reply to Professor R. Nielsen
%  Kjøbenhavn: P. G. Philipsens Forlag, 1868.   31 pp.
%
%  TRANSLATION (English) --- Hans Halvorson
%  Translated from transcription.tex (Danish), which is COMPLETE and
%  image-verified, printed pp. 3--31 + the author's Rettelse leaf.
%  Method: ../../../TRANSLATION-PLAYBOOK.md.
%
%  STATUS: in progress. Batch markers below, in reading order.
%
%  ---- WHAT THIS PAMPHLET IS ----
%  The last round of the exchange begun by Brøchner's _Problemet om Tro og
%  Viden_ (1868), to which Nielsen replied with _Hr. Professor Brøchners
%  philosophiske Kritik, gjennemseet af R. Nielsen_. Brøchner is answering
%  that piece. It is continuous polemical prose --- no chapters, no sections,
%  only three short centred rules --- and it is far more sardonic than
%  anything in _Problemet_. The English should keep the sarcasm rather than
%  flattening it into neutral exposition: the military metaphors (Angreb,
%  Tilbagetog, offensivt Stød, Baghold, Manøvre), the Olympian-thunderer
%  joke on p. 3, the cuttlefish on p. 7 and the Elverdans on p. 10 are the
%  argument's manner, not decoration.
%
%  ---- TERMINOLOGY: ALIGNED WITH ../problemet-tro-viden/translation.tex ----
%  This pamphlet stands beside its parent volume, which is already complete
%  in both languages. Where the two overlap, THIS FILE FOLLOWS THAT ONE.
%  Checked against it, occurrence counts in brackets:
%
%    (absolute) Uensartethed = (absolute) heterogeneity   [52 in the parent]
%       NB Brøchner spells it „Uensartethed“ with ONE e; Nielsen's own books
%       spell it „Ueensartethed“. Same term, same English.
%    Erkjendelse           = cognition        [367 in the parent --- NOT
%       "knowledge"; the Nielsen translations in this repo use "knowledge"
%       for Erkjendelse, and that is the one real collision between the two
%       authors' glossaries. Brøchner's own word wins in Brøchner's own book.]
%    Viden                 = knowledge
%    Videnskab             = science (Wissenschaft; never natural science)
%    Bevidsthed            = consciousness
%    Forestilling          = representation
%    Tro / Troen           = faith
%    Villie                = will;  Enkeltvillie = individual will
%    Underet               = the miracle;  Underbegreb = concept of miracle
%    Personlighed          = personality
%    Velleitet             = velleity (keep the term; it is Brøchner's jab)
%    Piecen                = the piece (Nielsen's pamphlet; Brøchner uses the
%                            word with visible distaste --- keep it, do not
%                            soften to "pamphlet")
%    Mening                = "opinion" where Brøchner is mocking Nielsen's
%                            refusal of the word Theorie; else "sense"
%
%  ---- CONVENTIONS ----
%  - Quotes. Danish „...“ -> English curly doubles ``...''. The book has THREE
%    quote patterns; guillemets «...» follow the FOUNT, not the sense, and are
%    used at exactly two places (p. 19 «lazzi», p. 21 «practica ... arithmetica»).
%    Those two keep guillemets in the English too.
%  - Emphasis. \emph{} in the Danish is letterspacing in the print; carry every
%    one over. Brøchner letterspaces heavily and pointedly --- the emphasis IS
%    the sarcasm in several places (e.g. „altsaa dog \emph{Theorie}!“).
%  - Latin/foreign in \textit{} stays \textit{} and stays untranslated
%    (calendæ græcæ, Nomina sunt odiosa!, instar omnium, lucus a non lucendo).
%    Where the print letterspaces a word INSIDE an antiqua phrase it is
%    encoded \emph{\textit{...}}; that nesting is reproduced here.
%  - Footnotes: \renewcommand{\thefootnote} gives *) **) ***) then arabic, and
%    \setcounter{footnote}{0} resets at each printed page, exactly as in the
%    Danish. 16 notes in the book. TWO of them RUN OVER a page break in the
%    print (p. 4's **) onto p. 5, p. 10's *) onto p. 11); each is set IN FULL
%    at its call in both files, with a comment at the following page marker.
%  - Page markers "% --- p. N ---" carried at the same joints.
%  - The old „ɔ:“ id-est mark -> "i.e." (two occurrences, pp. 8 and 13).
%  - Nielsen's and Brøchner's own work-titles stay in Danish inside references
%    (Om theor. og prakt. Erkj.; den gode Villie; Grundideernes Logik).
%
%  ---- DEFECTS: DANISH AS PRINTED, ENGLISH CORRECT ----
%  Same rule as the Nielsen volume in this repo. The Danish is the diplomatic
%  witness and carries every defect; the English reproduces the SENSE and logs
%  the divergence in a % note at the site. Danish „“ balance runs +1 ON PURPOSE
%  (p. 20, an outer quotation never closed); the English balances at 0.
%  Logged defects: p. 16 „poa“ for „paa“ · p. 20 „Adjertiv“ for „Adjectiv“ ·
%  p. 20 the unclosed quotation · p. 21 „geometica“ for „geometrica“ ·
%  p. 21 „samledes„ paa“ · p. 24 note **) closes without a full stop ·
%  p. 27 „Prof. N.'s vil kun“ where the sense wants „Prof. N. vil kun“.
%
%  ---- THE RETTELSE IS NOT APPLIED ----
%  Brøchner's own erratum leaf reads „Side 11, Lin. 2: videnskabelige, læs:
%  subjective.“ and it is RIGHT --- the sense wants „subjective“. Per the
%  repo's diplomatic stance the printed „videnskabelige“ stands at p. 11 in
%  the Danish, and the English likewise translates the PRINTED reading
%  ("scientific"), with a % note at the site. The Rettelse leaf itself is
%  translated at the end. Do not silently apply it.
% =====================================================================
\documentclass[12pt,a4paper]{book}
\usepackage[utf8]{inputenc}
\usepackage[T1]{fontenc}
\usepackage{libertinus}
\usepackage{libertinust1math}
\usepackage{textalpha}   % Greek (ἐνέργεια p. 14, ἄπειρον p. 22)
\usepackage{graphicx}
\usepackage[english]{babel}
\usepackage[protrusion=true,expansion=true]{microtype}
\usepackage{setspace}
\usepackage[margin=1.2in]{geometry}
\usepackage{enumitem}
\usepackage{fancyhdr}
\setlength{\headheight}{28pt}
\addtolength{\topmargin}{-13.5pt}
\usepackage{hyperref}

\hypersetup{
  pdftitle={A Reply to Professor R. Nielsen},
  pdfauthor={H. Brøchner},
  hidelinks
}

\renewcommand{\thefootnote}{%
  \ifcase\value{footnote}\or *)\or **)\or ***)\else \arabic{footnote})\fi}

\pagestyle{fancy}
\fancyhf{}
\fancyhead[LE,RO]{\thepage}
\fancyhead[RE]{\textit{A Reply to Professor R. Nielsen}}
\fancyhead[LO]{\textit{H. Brøchner}}

\setstretch{1.3}

\title{\textbf{A Reply\\[2pt] to\\[2pt] Professor R. Nielsen}}
\author{Dr.\ H.\ Brøchner\\[4pt]
  \small Professor of Philosophy\\[6pt]
  \small Translated from the Danish by Hans Halvorson}
\date{Copenhagen: P.\ G.\ Philipsen, 1868}

\begin{document}
\maketitle
\thispagestyle{empty}
\newpage

% =====================================================================
%  TEXT.  Printed pp. 3--31, continuous prose, no divisions.
%  Printed p. 1 is the title page and p. 2 the printer's imprint
%  (Bianco Lunos Bogtrykkeri ved F. S. Muhle); neither is transcribed as
%  body text -- the title page is represented by \maketitle above.
% =====================================================================

% --- p. 3 ---
% Printed p. 3 opens with an oversized ornamental Fraktur initial; the Fraktur
% sort serves both I and J and the word is „I den Efterskrift“, so it is
% resolved as I and set as ordinary type in both files.
\setcounter{footnote}{0}
In the Postscript that Professor Nielsen has appended to his ``popular''
lectures from Christiania, he concludes with a declaration, drawn up in
authoritative form and luxuriously furnished with the trinkets of street
language. Professor Nielsen will not reply, even should a fifth and a sixth
assailant join the four whom there is supposed to be no sense in answering ---
of whom, however, the two, the two ordained men, would perhaps, in accordance
with the Professor's old and ever-recurring love of theology, have been
answered.
Professor Nielsen withdraws majestically with that declaration, and now,
enthroned upon the Olympus of his two-peaked logic and shrouded in the
cloud-banks of his theory of faith, he will at most hurl --- not exactly the
Cloud-gatherer's lightning, but that species of projectile which drives off the
enemy although it can neither wound nor harm.

Then came the ``fifth assailant'' with his book, and Professor Nielsen made
haste, without waiting for the book's completion --- which was not, after all,
like Prof.\ N.'s philosophy of religion, first to await the completion of the
Logic of the Fundamental Ideas some
114½ years hence\footnote{% printed *)
  The calculation is made in Dr.\ Heegaard's book: Prof.\ R.\ N.'s Lære om Tro
  og Viden, p.\ 94 f.}
% „Dr.“ is ANTIQUA in the print here (and again on p. 11), as the arabic
% figures are and as the roman numerals „IX og X“ on p. 8 are. It is an
% abbreviation, not quoted foreign matter, so it is left plain rather than
% \textit{}-ed, in both files.
and so appear at the next \textit{calendæ græcæ}, but was to be given this
year; and he wrote a piece bearing the title: Hr.\ Professor Brøchners
philosophiske Kritik, \emph{gjennemseet} af R.\ Nielsen.
% EMPHASIS: within the cited title only „gjennemseet“ is letterspaced --- the
% surrounding words are set tight. The title is left in Danish, as a work
% Brøchner is citing; the letterspacing is carried on the same word.
% (In English: "Professor Brøchner's Philosophical Criticism, _Reviewed_ by
% R. Nielsen".)

% --- p. 4 ---
\setcounter{footnote}{0}
Now is this piece a reply, or merely a projectile, or is it a higher
dialectical unity of reply and non-reply, of reply and projectile? What follows
will give the reader the premises from which he must then draw the conclusion
himself.

In my criticism of Professor Nielsen's theory I had established that the
determination which was the theory's chief cornerstone --- the \emph{absolute
heterogeneity} which Prof.\ Nielsen accented so strongly, and which the echoes
repeated in chorus --- expressed something self-contradictory and impossible.

I had shown how this concept, which Prof.\ Nielsen could not hold fast and
would not give up, brought confusion in everywhere, and bred of itself a spawn
of contradictions, which I enumerated in my book, p.\ 213 f.\footnote{% printed *)
  Prof.\ N. concedes (the piece, p.\ 3, at the foot) that I have pointed out
  \emph{real} contradictions, but he avoids saying which. Might it perhaps be,
  among others, the \emph{absolute heterogeneity}?}

I had shown how the Nielsenian theory, while it undermines the ground of
knowledge, annihilates the ethical and corrupts the religious.

I had shown how this theory is in Prof.\ N. an empty thought-experiment,
without root in a personality, without truth in a life's actuality.

I had, with respect to the scientific work to which Prof.\ N. points back as
the foundation of his theory of faith, established its dualism, and pointed out
a thoroughgoing confusion in the fundamental concepts, a want of scientific
form. I had established the impossibility of the projected ``inverted''
science: the philosophy of religion (in the Nielsenian sense), the one by which
a distant future was to be made blessed.

To all this Prof.\ Nielsen has been able to answer \emph{nothing}\footnote{% printed **)
  % THIS NOTE RUNS OVER THE PAGE BREAK IN THE PRINT: it begins in the note area
  % of p. 4 and its last four lines stand, unmarked, at the head of the note
  % area of p. 5, above that page's own *). It is set here IN FULL at its call,
  % as in the Danish; see the % comment at the p. 5 page marker below.
  It is as though a teasing demon had, unbeknown to Prof.\ N., wheedled this
  admission out of him. In the very first lines of his book Prof.\ N. concedes
  that I have \emph{refuted} several other theories. The only theories I have
  treated besides Nielsen's are Martensen's and Kierkegaard's. But now I have
  made it clear that Nielsen's theory stands in such a relation to
  Kierkegaard's that the refutation of the latter is also the refutation of the
  former, while the relation certainly cannot be reversed. That admission
  therefore made it superfluous for Prof.\ N. to write 52 pages; he could have
  been content with 3 lines.}:
% --- p. 5 ---
% The note area of p. 5 OPENS with the unmarked last four lines of p. 4's **);
% they are set above, at the call on p. 4, and are NOT repeated here. The
% page's own note, *) „Nomina sunt odiosa!“, follows them on the leaf.
\setcounter{footnote}{0}%
it stands, now that his piece has appeared, as before, unrefuted. And yet, if
only the doctrine of the \emph{absolute heterogeneity} must be given up by
Prof.\ N., then the battle is lost for him; for this \emph{absolute}
heterogeneity, ``on account of which'' the reconciliation of faith and
knowledge in the same consciousness is to become possible, carries, as has been
conceded often enough, the \emph{whole doctrine}. \emph{But the word itself is
avoided in Prof.\ N.'s book.}
% Letterspacing here is fine-grained in the print: in „den absolute
% Uensartethed“ both words are spaced; in „denne absolute Uensartethed“ only
% „absolute“ is. The whole closing sentence is spaced, „Men“ included. The
% English reproduces the same distribution.

A \emph{reply} to the attack Prof.\ N. has thus not given; to that extent he is
in good accord with his declaration. But what, then, has been done in the piece
over the 52 pages?

There is, in the first place, an attempt to deny the correctness of my account
of Nielsen's theory --- which, moreover, I am not readily permitted to call a
theory. Well, about the name I will not quarrel; only I come to be a little
embarrassed for finding another designation. That I cannot call the thing which
must \emph{not} be called Prof.\ N.'s \emph{theory} Prof.\ N.'s
\emph{practice}, I have established to the point of evidence in my book
(p.\ 214 f.). But we might of course call it Prof.\ N.'s \emph{opinion}, for an
opinion is after all a ``\emph{neutral}'' expression; or else we might call it
Prof.\ N.'s \emph{mystery}, which might be fitting for the reason that he is
reluctant to come out with it, reluctant to ``waste a word'' on illuminating
it, but is glad to shroud it in the draperies of ambiguity. Still, let it stay
with ``the opinion''! So, then: I have not reported Prof.\ N.'s opinion
correctly; I have only, out of the Professor's occasional writings and out of
``\emph{another's}''\footnote{% printed *)
  \emph{\textit{Nomina}} \textit{sunt odiosa!}}
% „Nomina“ is letterspaced antiqua, „sunt odiosa!“ tight antiqua; encoded
% \emph{\textit{…}} and \textit{…}, as in the Danish. This is one of the
% antiqua phrases that carries NO guillemets.
accounts, put together for myself ``a something'' that is supposed to be that
opinion, and have not awaited the ``authentic and coherent'' presentation. That
this ``other's'' accounts are authentic in the \emph{strict} sense Prof.\ N.
will assuredly not deny: he knows that I
% --- p. 6 ---
\setcounter{footnote}{0}%
can furnish the proof of it. But perhaps the fault lies in their not being
``coherent''? They are given in systematic form, and their brevity is precisely
their merit; for the opinion shows itself in them, stripped of the padding of
phrases and the draperies of ambiguities, in its natural and naked shape, and
the ``crooked and skewed'' cannot be hidden.

What I am supposed to have misunderstood in Professor N.'s ``opinion'' is said
to be his \emph{concept of faith} and his conception of \emph{the miracle}.
With respect to this it has to a remarkable degree escaped Prof.\ N. what has
not escaped me: that velleities are not sufficient. I have, in the survey of
Prof.\ N.'s opinion, expressly brought out (p.\ 127) that he \emph{seeks} to
get cognition and thought into the will, to which faith leads back; but I have
in the criticism of the same opinion pointed out (e.g.\ pp.\ 196, 199) that
this striving comes to no more than a velleity, because the \emph{absolute
heterogeneity} makes that presence of cognition in the will an empty phrase, a
bare postulate\footnote{% printed *)
  Nor does it succeed any better for Prof.\ N. in the piece. Compare in
  particular the confused determinations on pp.\ 16 and 17.}.
% „Postulat“ is NOT letterspaced in the print (the sorts are tight);
% „U n d e r e t“ in the next sentence is. Same distribution here.
And it stands the same way with the conception of \emph{the miracle}. A run is
made at getting a lawfulness in: the miracle-working will of God is to be at
one and the same time one with the laws and raised above the laws, and by this
run one would so gladly keep off the awkward conflict, the conflict between
faith and knowledge. But the miracle-working God is for N., as I have shown,
``the God of arbitrariness,'' ``the God who \emph{creates out of
nothing}''\footnote{% printed **)
  \emph{Nielsen}: Om theor.\ og prakt.\ Erkj.\ pp.\ 70, 74 f.}, and this God
who creates out of nothing abolishes the laws just as arbitrarily as he has
given them; for they are for him only laws in that which is a nothing, and
therefore themselves a nothing.
% There is NO closing quotation mark after „skabende Gud“ in the print; the
% Fraktur OCR invents one. Nothing is supplied in either file.
It is only an empty word when it is said (p.\ 52): ``what is characteristic of
the miracle is not at all that a breach is made in the laws, but that a fact
\emph{is} immediately \emph{referred}\footnote{% printed ***)
  % the only ***) in the book
  Compare: ``\emph{that in which the believer sees the miracle!}'' (Nielsen:
  den gode Villie, p.\ 33).} to God's almighty
% --- p. 7 ---
will''; for precisely this \emph{immediate} referring to the \emph{almighty,
creating} will expresses a leaping over of the laws, a \emph{breach} of the
laws.
% The quotation „det for Miraklet Charakteristiske …“ opens on p. 6 and closes
% here, at the top of p. 7.

So, then: my conception of Prof.\ N.'s concept of faith and concept of miracle
is perfectly correct; and neither what the Professor \emph{willed} nor what
arrested this willing has been overlooked. And as regards the objection Prof.\
N. makes to the way in which I have understood his determinations about the
miracle's objective actuality --- that point, so highly vexing for Prof.\ N.,
which at once puts him in a passion --- I need only refer to what Prof.\ N. has
himself had the goodness to quote from my book. I have shown how Prof.\ N.'s
presuppositions inevitably had to drive him to the attempt, in order to
forestall the conflict between faith and knowledge, to \emph{double} actuality,
so that faith and knowledge each got \emph{its own} actuality. But thereby
Prof.\ N. came into the unpleasant situation, over against believers, of having
to take away the objective, outward actuality which they require for faith's
object, and which they cannot give up if they will not give up religion's
historical foundation. Prof.\ N. thus stood between Scylla and Charybdis. If
objective actuality was held fast ``where it really constitutes an
indispensable moment in the phenomenon of faith,'' then the conflict with
knowledge became inevitable; and if objective actuality --- for example of
Christ's birth, death and resurrection --- was transformed into a merely inner
actuality, then faith protested. For this faith knows very well ``how to make a
distinction between outer and inner actuality''; it is not confused by phrases,
and I have surely shown plainly enough what a world of difference there is in
this respect too between \emph{real faith} --- the faith that is alive in works
--- and \emph{the faith of imagination} --- the faith that is lively in
declamations. Prof.\ N. stood and stands, then, between two dangers, and can
save himself only in the way the cuttlefish saves itself: by making it black
around him. I shall follow Professor N.'s example, and refer his readers to the
Christiania
% --- p. 8 ---
% „Christianiaforelæsningerne“ is broken by the compositor at the foot of p. 7
% and rejoined with \-% in the Danish so the transcription reflows.
lectures IX and X, and then only ask them, with a little reflection, to read
through the reply to Prof.\ Caspari in the Postscript; they will then
understand how the device is applied: ``to talk black.''
% The roman numerals „IX og X“ are ANTIQUA in the print, as all figures and
% roman numerals in this book are; set plain in both files, not \textit{}.

Prof.\ N.'s objection to my conception of his opinion we may then regard as
dispatched. But once it is removed, the Professor has a new turn ready to hand,
and one over which he seems to be exceedingly glad: by my criticism of ``the
opinion'' I am supposed to have achieved precisely what the crafty strategist
against whom I was fighting himself wished --- to confirm what Prof.\ N.
inculcates again and again: that the doctrine of faith cannot possibly stand
before science's tribunal, and that it therefore neither can nor will submit to
science's decision. In my rashness I have thus fallen into the craftily laid
ambush, and my crafty opponent seems exceedingly glad. Yet this gladness might
perhaps not be quite so reliable a thing; for Prof.\ N. sees himself compelled
to perform some little feats of dexterity in order to convince others that it
is well founded. What I have directed my criticism against is the opinion about
\emph{the relation of faith and knowledge}, a relation which, by Prof.\ N.'s
own determination, knowledge recognizes, in that it mirrors the opposition
while it comprehends its own essence. My \emph{criticism of knowledge} is thus
directed against \emph{determinations of knowledge}; my tactics are thus, on
the presuppositions, entirely correct. But in order that they shall not be so,
Prof.\ N. first lets the thing I was not permitted to call a theory become
``the \emph{theory of faith} set up by me (i.e.\ Prof.\ N.)'' --- so it is a
\emph{theory} after all! --- and then, by a quick little maneuver, lets
\emph{the theory of faith} become \emph{the doctrine of faith}.
% ɔ: -- the id-est mark (reversed c + colon) in the Danish, rendered "i.e."
% here. The Fraktur OCR reads it „o:“ and the ABBYY layer „໑:“.
And then of course I am convinced of having let myself be caught. But even if
we now let this substitution stand, I have in my book established precisely how
the judgment upon everything that becomes an expression of the human essence
belongs in the last instance to knowledge, in virtue of knowledge's relation to
the essence; I have not postulated it, as N. asserts (the piece, p.\ 8), but
have furnished the ``fundamental
% --- p. 9 ---
% „Grundbeviset“ is broken at the foot of p. 8 and rejoined in the Danish.
\setcounter{footnote}{0}%
proof'' in both \emph{direct} and \emph{indirect} form, and established the
validity of the proof from knowledge\footnote{% printed *)
  See in my book the developments summed up on p.\ 210 f., cf.\ p.\ 169.}.

Prof.\ N.'s gladness over the crafty turn was thus short-lived; he must now set
about carrying out the retreat, and it is then to be covered only by a little
offensive movement. By that movement it first becomes quite interesting that
Prof.\ N. comes to commit a tactical error of the same kind as the one he
wished above to foist upon me. For even if the ``offensive thrust'' had turned
out ever so successfully, even if he had annihilated everything I had given
through my criticism by way of positive contributions to a psychological and
metaphysical theory, Prof.\ N. would still be no nearer. He would have stood
face to face with the impossibility of holding fast the ``\emph{absolute
heterogeneity},'' that determination which carries Prof.\ N.'s whole theory ---
now we may of course call it \emph{theory} again! --- and with which it stands
and falls. And when he then gave it up --- which he was perhaps by now not so
disinclined to do, if it could be done in a tactful way --- and got himself a
new theory put together which was to lead to the same goal as the old, he would
stand face to face with the insoluble problem of how, when knowledge's
normative validity for the whole of personality is denied, there can be found a
norm which on the field of the ethical and the religious makes a separation
between the true and the untrue, between personality's true and essential
demand and the individual's untrue striving\footnote{% printed **)
  Cf.\ my book, e.g.\ pp.\ 169, 203.}. I have in my book pointed out the
impossibility of it: Prof.\ N. will be as little able as anyone else to
establish its possibility.

However Prof.\ N.'s attack had turned out, then, his theory would have been
just as little saved; my book's negative result would stand unshaken. But
Prof.\ N.'s attack is so little dangerous that it can be thrown back with
slight effort --- indeed it has essentially already been dispatched in the very
book against which it is directed.

% --- p. 10 ---
% p. 10 carries only SEVEN lines of body text; the rest of the leaf (some 33 of
% its 40 lines) is the single note *), which then runs over onto p. 11.
\setcounter{footnote}{0}
Of the attack itself it will of course be only the chief points that I come to
touch on; to go into the individual little turns, little witticisms and little
inferences with large paralogisms will be quite superfluous in the face of any
reader who is not either by nature very muddled or has become so through
incautious enjoyment of the sort of logic that is served up in the ``60
sheets''\footnote{% printed *)
  % THIS NOTE RUNS OVER THE PAGE BREAK IN THE PRINT: it fills the note area of
  % p. 10 and its last two lines stand unmarked in the note area of p. 11,
  % which carries no note of its own. It is set here IN FULL at its call, as in
  % the Danish; see the % comment at the p. 11 page marker below.
  Thus I think I may confidently leave it to any reader who does not suffer
  from the above-mentioned natural or acquired weakness to come to terms for
  himself with what Prof.\ N. develops on p.\ 8 f. about \emph{existence and
  theory} and on p.\ 20 f. about \emph{the will's relation to the self and its
  self-determination}. One need only note what Prof.\ N. cautiously slips past,
  and what he, not very deftly, slips in, in order to get one's bearings. Just
  as calmly do I think I may leave to him the criticism of the astonishing
  consequences that Prof.\ N. (p.\ 24 f.) extracts from what I have said about
  the unity of a \emph{definite}, particular form of theoretical thinking and
  that form of thinking which would be, on the presupposition, \emph{absolutely
  heterogeneous} with all theory --- ``practical thinking.''
  % „ab=/solut uensartede“ is broken across a line in the print and the
  % letterspacing carries across the break; rejoined in the Danish.
  Nor would it be beyond him to correct for himself what Prof.\ N. infers from
  the fact that one personality can never become completely transparent to
  \emph{another} personality, or to see what difference I make between Prof.\
  N.'s ``absolutely heterogeneous'' faith and faith in the sense of the
  religions. Finally, it would surely not be difficult for him either, by the
  ``example \textit{instar omnium}'' of my ``empty plays on words,'' over which
  Prof.\ N. is so delighted (p.\ 50), to make the discovery that the emphasis
  falls on this: that the mystery which, in Prof.\ N.'s opinion, was to
  \emph{prevent} the mutual limitation of faith and knowledge comes precisely
  to \emph{express} the limitation; and that the ``fine'' point about the
  ``mystery'' \emph{not} becoming something mysterious is that a secret is no
  longer a secret once one knows \emph{what its content is}. But precisely that
  is known when faith and knowledge are each other's mystery, and are present
  in one and the same consciousness. Here, however, the ``absolute
  heterogeneity'' --- this great and fatal discovery of Prof.\ N.'s --- has
  played his ``good will'' the trick that the consciousness which is to hold
  both faith and knowledge really is ``split.'' Let one read in ``den gode
  Villie'' (p.\ 33): ``when the believer and the knower'' (who, be it noted,
  are the \emph{same} individual) ``\emph{both} at once direct their gaze
  toward the mystery, the strife is composed and the extremes are reconciled''!
  --- So, just as in the Elf-Dance the One became Three, here the One becomes
  Two; but the Two will, ``on account of'' the absolute heterogeneity,
  unfortunately not become One again.}.
% --- p. 11 ---
% The note area of p. 11 holds ONLY the unmarked last two lines of p. 10's *).
% They are set above, at the call on p. 10, and are NOT repeated here; p. 11
% has no note of its own and no \setcounter is needed.
My first sin is said to be that I confound the general's scientific
% RETTELSE. Brøchner's own erratum leaf reads „Side 11, Lin. 2: videnskabelige,
% læs: subjective.“ and it is right: the sense wants "subjective validity with
% the objective". Per the repo's diplomatic stance the PRINTED reading stands
% in the Danish and is translated here; the Rettelse is set at the end of the
% edition and is NOT applied.
validity with the objective, and that to such a degree that I do not see myself
able to make a distinction between grounds of science and grounds of
conscience, between a cognition of the law of duty and a cognition of the law
of gravitation. And this sin of mine is the gravest, which is why Prof.\ N.
ends his piece by calling me to account for it. But here I am so fortunate as
already to have rendered the account, which may be inspected in my book,
namely pp.\ 141---148. Only with respect to my lack of ability to make a
distinction between a cognition of the law of duty and of the law of
gravitation must I point out that this juxtaposition is Prof.\ N.'s, not mine.
I have juxtaposed the theoretical cognition of the human being's essence and
the cognition of the duty that is determined by the essence, which becomes
manifest to itself in that theoretical cognition; and on account of duty's
relation to the essence I have held fast the unity of theoretical cognition and
\emph{cognition of duty}, the significance of \emph{theoretical} cognition as a
necessary condition of the truth of action. If Prof.\ N. in his ``development
of the case'' had not used the procedure which Dr.\ Heegaard so sharply and
aptly characterizes in his little book, p.\ 29 f., but had, instead of the
objective cognition of the law of gravitation, which ``does not concern the
world of personality at all, but only that of corporeality,'' put an objective
cognition of what falls within the psychological --- of the human self itself,
of human personality and its fundamental functions --- then the one who
cognizes would not have been ``so expressly called upon to disregard all
personality,'' and then the ``development of the case'' would perhaps have led
to another result.
% „Dr.“ is antiqua here as on p. 3; left plain, see the comment at p. 3.

My second sin is said to be that by a mistake I put the general concept of the
individual will in the place of the will that is struggling for its
independence. May I put a question to Prof.\
% --- p. 12 ---
N.? When it is to be determined what the individual will \emph{as individual}
does in determining itself to realize in action what has been cognized, how
shall it then be possible to \emph{speak} of this act of will, if one does not
speak in \emph{concepts} which, as the expression of the individual will's
essence, express the individual will's law?
% PRINTER'S DEFECT, p. 12: a solid round blot of ink stands in the word space
% between „Væsen“ and „ud-/trykker“ in the Danish. It is foul matter, not a
% sort; the two words are transcribed there with an ordinary space, and
% nothing is reproduced here.
This is not to identify the general concept of the individual will with the
individual's will, but to posit the latter as that in which the former acquires
actuality, and which therefore has its valid law in it. When, on the other
hand, Prof.\ N. speaks from his standpoint of the concept of will, there is
this peculiarity about it, that this concept of will does \emph{not} become a
concept of the will, and is thus called a concept of will in the same way as
\textit{lucus a} \emph{\textit{non}} \textit{lucendo}.
% Antiqua, with ONLY „non“ letterspaced: the print gives „lucus a n o n
% lucendo.“ Encoded \textit{…} for the tight antiqua and \emph{\textit{…}} for
% the spaced word, as in the Danish. No guillemets here.

For the will, on Prof.\ N.'s conception, becomes --- however much he struggles
against it --- that which, by its absolute heterogeneity with knowledge, cannot
find an expression in the concept. One would then in any case have to let the
common determination \emph{concept} comprise two absolutely heterogeneous
subspecies, as the subspecies of Prof.\ N.'s concept of creation and concept of
miracle do (cf.\ my book, p.\ 206 f.); but such a relation is an impossibility,
and it becomes only an empty play with words to designate the absolutely
heterogeneous by the same determination. Should Prof.\ N. wish to call his
concept of will a ``concept of mirroring,'' I need only refer him to my
elucidation of this ``category'' --- mirroring (in my book, p.\ 148 ff.).

So we come to my third and last sin: my lamentable confusedness in the
determination of the energy of selfhood and of the relation of knowledge and
will to the fundamental concept of the essence. Here, however, it has happened
to Prof.\ N. that he has in a very peculiar way misunderstood certain
fundamental determinations I have given --- determinations which I have never
had difficulty in getting the youngest students to understand, and which the
Professor too would surely easily have grasped correctly, had he not so
immersed himself in his own terminology that he rightly understands only what
is expressed in it.

I have, in the psychological developments in my book, brought out
essence-universality as that which, in contrast
% --- p. 13 ---
to the particularity of the animal essence --- a particularity that betrays
itself in the life of instinct --- characterizes the human essence.
% PRINTER'S DEFECT, p. 13: a thin dotted vertical stroke stands in the word
% space before the second „Particularitet“ in the Danish. Foul matter, not a
% sort; transcribed there as an ordinary word space, nothing reproduced here.
In this determination of essence, which encloses the determination of inner
infinity, I have sought the possibility of self-consciousness --- the
possibility, that is, of that relation of the self to itself, that being-for-self
of the self which, in contrast to the animal self's being-for-self in
self-sensation (in which the universal is indeed present, but always as bound to
the \emph{individual} sensation), may naturally be said to have the liberated
--- i.e.\ the form of universality liberated from that boundness; and in virtue
of cognition's relation to essence-universality (p.\ 137) I have ascribed to it
the ideal primacy in the human spirit.
% ɔ: -- the id-est mark (reversed c + colon) in the Danish; rendered "i.e." here.
% RECON.md §8 records this as the one confirmed instance in the book.

Prof.\ N. has now in a peculiar way misunderstood the significance of the
determination \emph{essence-universality} and of the designation \emph{the
liberated universality}, and in virtue of these misunderstandings he draws a
mass of inferences of a quite singular kind (p.\ 26 ff.), through which
absurdities are demonstrated against me. But as soon as the misunderstandings
are removed, they take the absurdities with them.

Then I have, in my book p.\ 135, determined the relation of cognition and will
to ``the energy of the essence revealed as unity in self-consciousness,'' and
for the determinations I have given there I must hear a great deal of ill from
my esteemed critic. First, I am supposed to have borrowed my ``catchwords'' ---
selfhood and energy of selfhood --- from Prof.\ N.'s philosophy. That is a
little surprising to me. Not only have I used these categories in psychological
lectures delivered before any of those writings of Prof.\ N. had appeared in
which he has \emph{touched} psychological questions --- anything coherently
psychological Prof.\ N. has never given. But I also stood in the naive belief
that selfhood --- the Danish word for subjectivity --- was a current category in
post-Kantian philosophy: a
% „courant“ is Fraktur in the print, not antiqua; so is „à la“ on p. 18. Only
% Latin and Greek are set in antiqua in this book (RECON.md §5).
category which neither Prof.\ N. nor I had invented or needed to invent. And as
for energy of selfhood --- an expression which I do not, for that matter,
remember having come upon in Nielsen, who in any case, in accordance with the
absolute heterogeneity, would not be able to make much
% --- p. 14 ---
of it --- it is the determination of selfhood combined with the old determination
of ἐνέργεια, which we owe to psychology's founder, Aristotle. It is from
Aristotle's
% Greek, antiqua italic in the print: smooth breathing over the first epsilon,
% acute over the second -- ἐνέργεια. Copied verbatim.
writings that I have fetched this determination, not from Prof.\ Nielsen's; and
although there sometimes betrays itself in the latter a striking unconscious
inclination on their author's part to annex what belongs to older philosophers,
I can hardly
% PRINTER'S DEFECT, p. 14: a short high horizontal stroke stands in the word
% space between „saa“ and „kan“ in the Danish; ink on spacing material, not a
% sort. Transcribed there as an ordinary word space; nothing reproduced here.
imagine that Prof.\ N. would fall into the self-delusion of supposing that he,
perhaps in his youth, at one of his many earlier stages of development, had
written the psychology which has since for some considerable time gone under
the name of a certain Aristotle. Of Prof.\ N., as I say, I cannot imagine it;
but for his
% PRINTER'S DEFECT, p. 14: a round ink dot at mid-height stands in the word
% space before „Navn“ in the Danish. Same class as the two marks above;
% transcribed there as an ordinary word space.
adherents --- or adherent, for now one ought perhaps to speak in the singular
--- I dare not answer; he has after all already discovered recognizable traces
of his master in Norse mythology.

Those ``catchwords,'' which I have thus \emph{not} borrowed from Prof.\ N., I am
then supposed to have handled badly in determining them more closely. I am
supposed to have got out a peculiar trinity of energies of selfhood, whose
confusion in Prof.\ N.'s description (p.\ 38 ff.) looks almost as frightful as
the confusion among his fundamental ideas. But a couple of short remarks may
perhaps show that the whole of Prof.\ N.'s consequence-drawing rests on
misunderstandings. What I have said is the following: the fundamental
determination of the human essence is energy of selfhood; in the unity of
self-consciousness the unity of the essence reveals itself, and for the energy
of this undivided essence cognition and will are primitive expressions,
primitive forms of realization. The general fundamental determination of the
energy of selfhood thus receives its particular expressions in the forms of the
energy of cognition and the energy of will; in both these primitive forms the
essence expresses itself in such a way that its unity is preserved --- not, as
Prof.\ N. will have it, \emph{split apart} in absolute heterogeneity --- and in
both forms the content is therefore \emph{essentially} the same; for the content
is a determination of the essence, and the forms' abstract opposition is
abolished in that the determinations which express it point to one another and
develop out of one another (cf.\ my
% --- p. 15 ---
\setcounter{footnote}{0}
book, p.\ 134 f.). The fundamental forms must, precisely as fundamental forms,
reflect one another's formal determinacies.

When one compares these propositions with what Prof.\ N. has written with
visible satisfaction and innumerable repetitions in his piece, p.\ 36 ff., one
will see that the Professor might, with a little less heat and a little more
reflection, have spared himself a number of un-shrewdnesses.

Should the reader, with respect to some single point --- for example with
respect to the relation of cognition and will at the imperfect stages of
development and at their culminating stage, or to the \emph{conversion} of the
content of cognition into the content of will and conversely --- desire closer
information, I refer him to my book (p.\ 136 ff.), and in particular to what
Prof.\ N. has left out in his citations\footnote{% printed *)
  Prof.\ N. cites on the whole somewhat carelessly. Thus the citation found in
  his piece, p.\ 22, lines 6---11, is given in an entirely senseless form, and
  the corruption is moreover repeated within quotation marks further down the
  same page.}.

Beside the striking misunderstandings of which Prof.\ N. has made himself guilty
with respect to the points adduced, the accusations of misunderstanding that he
directs against me acquire a peculiar light. I have already shown how matters
stood with my supposed misunderstanding of Prof.\ N.'s concept of faith and
concept of miracle; and with a couple of words I shall be able to show how other
supposed misunderstandings dissolve into perfectly correct understanding. I am,
according to the piece p.\ 32, supposed in my conception of the will to have
made myself guilty of so remarkable a misunderstanding that no one who ``might
reasonably be assumed to be able to read'' would be capable of it. The
misunderstanding is supposed to consist in my having traced Prof.\ N.'s concept
of will (\textit{sit venia verbo!}) back to the
% ANTIQUA: „sit venia verbo!“ is upright antiqua standing against the Fraktur,
% like the other Latin in this book, and carries no guillemets.
``immediate expression of activity of natural determinacy in the drive, the
sensuously egoistic craving,'' and against this Prof.\ N. sets himself with
hands and feet. But, as I have shown above, velleities are of no help, and with
the absolute heterogeneity between knowledge and will as
% --- p. 16 ---
% „Forudsætning“ is broken by the compositor at the foot of p. 15 and rejoined
% in the Danish with \-%.
presupposition, the will, in its separation from the universal of knowledge,
really cannot become anything other than what I have let it become. It truly
does not help, in order to save a cognition and a thinking, and thereby the
determination of the universal, for this will, to posit two \emph{absolutely
heterogeneous} sorts of cognition and thinking; for we have already seen that
that is only to make the determination \emph{cognition} and \emph{thinking} an
empty word. So here there is no misunderstanding on my side; it is only that
Prof.\ Nielsen has not understood to what consequences the absolute
heterogeneity --- this absolute heterogeneity, about which he is now so silent
--- must inevitably lead.

In another place (p.\ 12) Prof.\ N. informs his readers that I have made myself
guilty of ``more than one distortion'' of his words, of ``a thoroughgoing
misunderstanding of the whole matter under discussion.'' This matter is the
relation between subjectivity and objectivity. The Professor takes on very
strongly and uses several pages to instruct and chastise me. And what is the
whole of it? A misunderstanding --- not on my side, but on Prof.\ N.'s. I had,
as a kind of introduction to the criticism of Prof.\ N.'s theory, illuminated
the confusion in his conception of some of the determinations by which he
expressed the fundamental opposition: the determinations existence and theory,
subjectivity and objectivity. I wished partly to make clear why in the criticism
I kept to another expression for the opposition, partly to illuminate from a new
side the relation between Kierkegaard and R. Nielsen by showing what fate the
fine and acute thinker's clear categories underwent on coming into Prof.\ N.'s
hands. I then pointed out the confusion and ambiguity in Prof.\ N.'s
determinations, and the sophistical use he made of them. With respect to the
opposition between subjectivity and objectivity, as Prof.\ N. determined it, I
drew attention among other things to two things. The one
% PRINTER'S DEFECT, p. 16: the Danish print reads „poa“ for „paa“, kept as
% printed there. The English has the correct word.
was that, while the abstract opposition is definitively held fast, it is
arbitrarily ignored that at the various stages of subjectivity the subject's
determinacy ---
% --- p. 17 ---
\setcounter{footnote}{0}
as appears especially clearly at the stage designated as that of infinite
subjectivity --- is precisely in virtue of what has been determined as
existence's objective principle, and that it expresses this. The other was that
it is overlooked that subjective truth is precisely the appropriated truth, and,
in the appropriation, the truth essentially equally objective with itself. By
not understanding that the last passage expresses the relation generally, and
by referring it, in a highly arbitrary manner, exclusively to a single one of
the stages named, Prof.\ N. gets out that I have entirely misunderstood his
words about this, and that I have entangled myself in ghastly contradictions. As
soon as my words are understood correctly, one sees that they presuppose no
misunderstanding of Prof.\ N., and that they are perfectly correct. It shows
itself especially clearly at the standpoint of Stoicism and of Hedonism that the
subject's determinacy expresses the determinacy of existence's objective
principle. This needs no explanation for anyone who knows even a little of the
history of philosophy, and neither does it need an explanation that existence's
principle is designated as existence's objective principle. And when the general
proposition about the relation between subjective and objective truth is then
applied to these two systems, the consequence is that the relative truth
possessed by the individual one-sided system --- which in its one-sidedness
nevertheless encloses something valid --- is the measure of the subjective truth
in him who practically appropriates the system's objective truth\footnote{% printed *)
  To Prof.\ N.'s droll argument from the fact that a Stoic can ``put himself
  just as securely into'' the hedonic theory as a hedonist and yet not become a
  hedonist, I shall only direct the reader's attention: it needs no commentary!}.

Thus, then, matters stand with my misunderstandings, and such is the character
of the attack by which Prof.\ N., as by a rearguard action, seeks to cover his
actual retreat. But Prof.\ N. does not make this retreat willingly; he has grown
angry, exceedingly angry, heated, exceedingly heated. And this anger and this
heat he must give vent to in
% The signature numeral „2“ stands at the foot of p. 17 and is not transcribed.
% --- p. 18 ---
\setcounter{footnote}{0}
many ways. Now I must hear ill of my ``noticeable self-regard''; now because on
``an airy substratum of essence-empty abstractions, skewed deductions,
smeared-out\footnote{% printed *)
  A hard accusation in the mouth of the man who has written the second part of
  Grundideernes Logik!} turns of phrase and short-winded remarks'' I have
``stilted myself up into a philosophical authority, a high priest of
knowledge.'' Now my criticism is honored with the title of ``empty criticism of
words,'' and remarks à la Per Degn are ascribed to me. Then I am informed that I
lack respect for science and any sense for it, and then at last Prof.\ N.'s
benevolent mood culminates, in that he undertakes to seek motives for my
conduct. About the last point I shall presently have a serious word with the
Professor; one thing and another of the rest I shall here illuminate briefly ---
some perhaps slightly unwelcome side-lights will fall on my opponent at the same
time.

Prof.\ N. would reduce my criticism to a criticism of words, lacking ``all
development of the case.'' This turn Prof.\ N. has borrowed from that one
adherent, who has however already used it so often --- first against Cand.\
Brandes, then against Dr.\ Heegaard --- that it might now be a little worn out.
The ``example \textit{instar omnium}'' of
% ANTIQUA inside the Fraktur quotation marks: the „…“ are Fraktur, the Latin is
% antiqua and carries no guillemets. Same phrase as p. 10.
``empty play on words'' which Prof.\ N. adduces against me I have already
illuminated above (p.\ 10), and it turned out to be a good deal more than a play
on words. Conversely I have illuminated Prof.\ N.'s ``development of the case''
of the relation between the cognition of
% --- p. 19 ---
the law of gravitation and of the law of duty, and it turned out that it did not
exactly develop the case very far. It would be letting Prof.\ N. off far too
lightly if, every time his conceptual developments were shown to be false and
confused, he could simply cry ``criticism of words'' and then ``generally, in
all generality'' insist on development of the case. Thoughts are had only in
words, and they would have to be very vague and loose thoughts that could not
bear to be held fast in the verbal form in which they present themselves. And
development of the case is not exactly Prof.\ N.'s strong side. Only let one not
be dazzled by Prof.\ N.'s
habit of taking up, in his writings, great masses of material from the
heterogeneous sciences with which he dabbles; and let one be still less dazzled
by his speaking pompously of ``combustion of material in the fire-air of
dialectic'' and so on. The ``combustion of material'' is very slight, and on
somewhat closer inspection it will appear that in the great majority of cases
there is no more organic connection between Prof.\ N.'s conceptual developments
and the material taken in than there is in a zoological museum between the
animal hides and the tow with which they are stuffed. But there are certainly
those who have had so much tow put in their eyes that they cannot see it!

Then Prof.\ N. makes me into a second Per Degn! Well, well! Rasmus Berg could
after all make Per Degn into a cock; why then should Rasmus Nielsen not be able
to make me into Per Degn? I should not, to be sure, have wasted a word on it, had
not Prof.\ N., in order to get an opportunity to place his witticism, dealt with
citations in a somewhat peculiar manner. If Prof.\ N., instead of beginning the
citation from his own book with the words ``If it stands fast that faith etc.,''
had taken along the nonsense that precedes it, and if he, instead of beginning
the citation from my book with ``The lively analogies etc.,'' had taken along
what precedes (p.\ 187), in which that nonsense is pointed out, he would have
oriented the reader considerably better. And had he cited the note I appended in
full, the reader would at once have been able to see from whom the Per-Degnish
enlightenment comes.
% The print sets „Per = Degnske“ with the Fraktur hyphen and a word space on
% either side of it; transcribed „Per-Degnske“ in the Danish.

I had said in the note, on the occasion of some \textit{«lazzi»} of Prof.\ N.'s:
% ANTIQUA with ANTIQUA GUILLEMETS. The guillemets follow the fount, not the
% sense: most of the antiqua Latin in this book carries none (RECON.md §3).
% Both marks are inside \textit{} because both are antiqua. Kept in the English.
``Nielsen's examples with the `mathematical friends' and the `friendly
mathematics' elucidate only the not very interesting point that one cannot
without further ado, because an adjective can be predicated of a substantive,
predicate an adjective corresponding to the substantive of a substantive
corresponding to the adjective; that one cannot, for example, because one can
speak of an insipid witticism, speak of a witty insipidity. That
% Single-level nesting, as printed: the outer „…“ opens here and closes on p. 20
% after „gjennem dem.“, and the inner pairs round „mathematiske Venner“ and
% „venskabelige Mathematik“ are ordinary „…“, not the doubled „„…““ of p. 20.
% The English renders the inner pairs as `…' and the outer as ``…''.
% The signature numeral „2*“ stands at the foot of p. 19 and is not transcribed.
% --- p. 20 ---
there can be ethical and unethical astronomers elucidates only that different
activities of spirit can have the same subject without being capable of being
made predicates of one another (as, for example, within the sphere of knowledge,
astronomy and medicine, since there are indeed medically skilled astronomers but
no medical-scientific astronomy); but it proves nothing with respect to an
absolute heterogeneity, or to these expressions of spirit being predicable of the
same subject without their essentially going back to the same principle, or
without the same undivided essence asserting itself through them.'' In Prof.\ N.
one now reads (p.\ 50, note) the following: ``Instead of elucidating something
about the case, Prof.\ Brøchner has, perhaps on personal grounds, preferred to
elucidate something about the words.
% „Brøchner“ carries the ø slash here; BOTH OCR witnesses drop it.
The elucidation is `that one cannot without further ado, because an adjective can
% PRINTER'S DEFECT, p. 20, inside the quotation from Nielsen: the Danish print
% reads „Adjertiv“ for „Adjectiv“, kept as printed there. The English has the
% correct word.
be predicated of a substantive, therefore, conversely, predicate an adjective
corresponding to the substantive of a substantive corresponding to the
adjective' --- an elucidation which, as far as the case goes, stands about on a
level with Per Degn's elucidation of Imprimatur:
% „Imprimatur“ and the grammarians' terms are FRAKTUR in the print, not antiqua;
% no \textit{} in either file.
`All the words that can be named are referred to 8 things, which are Nomen,
Pronomen, Verbum, Participium, Conjugatio, Declinatio, Interjectio.'\,''
% DEFECT / QUOTE BALANCE, p. 20: in the Danish the doubled inner quotation closes
% here and the OUTER quotation --- opened at „Istedenfor --- is NEVER CLOSED in
% the print. Transcribed as printed there, which is the whole of that file's
% deliberate +1. The English closes the outer quotation here, where the sense
% ends, and therefore balances. The two doubled pairs on this page are the only
% „„…““ in the book; they are rendered `…' inside ``…''.
This elucidation is supposed to be mine; but from the citation of my book one can
see that the elucidation is precisely Prof.\ Nielsen's. I must therefore concede
him the honor of having spoken like Per Degn; but I appreciate his good will, and
shall in return presently be of service to him with a citation from the same play
of Holberg's, to which I can imagine Prof.\ N. is drawn by a natural sympathy.
But since I have come upon Prof.\ N.'s note, I shall permit myself to add a small
remark. Prof.\ N. has grown very angry that I have spoken of his attempts in the
burlesque, and he speaks somewhat enigmatically of ``decency'' and of ``prim,
hysterical, cantankerous old schoolmistresses,'' and of there being no ``impure
constituents'' in his examples --- words which he seems to ascribe to me as mine,
since he adduces them with quotation marks, although I have nowhere used them.
The Professor seems
% --- p. 21 ---
to have grown a little uncertain of his aesthetic determinations, and that may
easily happen to a man who has to study \textit{«practica, didactica, tactica,
homiletica, exegetica, ethica, rhetorica, oratoria, metaphysica, chiromantica,
necromantica, logica, talismanica, juridica, parasitica, politica, astronomica,
geometrica, arithmetica»} and much more --- that one thing and another may slip
his memory.
% ANTIQUA with antiqua guillemets, the whole list. PRINTER'S DEFECT in it: the
% Danish print reads „geometica“ for „geometrica“ (Holberg, Erasmus Montanus),
% broken across the line with the ANTIQUA hyphen and rejoined there. Kept as
% printed in the Danish; the English gives the correct „geometrica“.
When aesthetics comes round again among the Professor's many studies, it will be
advisable for him to refresh the ``conceptual development'' of the burlesque; the
``development of the case'' he will then be able to find in his own writings.

So we come to the Professor's accusation against me of a lack of respect for
science. This lack is established by a very remarkable proof, whose validity
rests in turn on one still more remarkable. The proof, the incontrovertible
proof, of my lack of respect for science is that I have used no more than 6 pages
to criticize the Logic of the Fundamental Ideas --- the epoch-making Logic of the
Fundamental Ideas! I stood, and still stand, in the belief that it was perfectly
sufficient if I briefly and clearly pointed out what an insoluble tangle there is
in ``the Fundamental Ideas'' Logic in the fundamental ideas --- precisely, that
is, in what was to hold the work together --- and if in a couple of examples,
taken from points of quite decisive significance for everything of the book that
has appeared, I showed the sleights by which the dialectic of word-play arrived
at its results. But in Prof.\ N.'s opinion that is by no means enough, and the
proof that it is not is supposed to lie in the fact that the Logic of the
Fundamental Ideas is ``no less than 60 sheets!'' Striking, irrefutable proof!
Remarkable scientific measure! Why does the Professor not rather state the work's
mass in pounds? Scientifically that measure would be just as good, and
practically it would be far more useful --- especially when the Logic of the
Fundamental Ideas, having ceased to be a compulsory article of use for the
younger students, is gathered ``in a place of its own outside literature'' with
% PRINTER'S DEFECT, p. 21: in the Danish the opening „ is set tight against
% „samledes“, with the word space AFTER it instead of before --- „samledes„ paa“.
% Kept as printed there; the English spaces it normally.
its elder siblings: the theological ``propædeutiske Logik'' of 1845 and Prof.\
N.'s logical firstborn, the poor
% --- p. 22 ---
% „stakkels“ is broken by the compositor at the foot of p. 21 and rejoined in
% the Danish with \-%.
little acephalous thing, the system ``without head or tail,'' the speculative
Logic of 1841, which ``systematically'' begins with § 11 and just as
``systematically'' ends in the middle of a sentence. The \textit{petitio
principii} of which Prof.\ N. is constantly guilty
% ANTIQUA, no guillemets. „§ 11“ is set with a space and without a full stop
% after the section sign here; p. 26 has both „§ 35“ and „§. 11“ --- the
% inconsistency is the book's.
with respect to the Logic of the Fundamental Ideas is that it is worth the
trouble at all to occupy oneself with this fragment of a logical (anti-logical?)
repertory. Had Prof.\ N. not, immersed as he is in his criticism of faith and his
criticism of knowledge, quite forgotten that there is a third and very useful kind
of criticism --- namely \emph{self}-criticism --- he would perhaps have gained a
clearer view of his own
% The compositor letterspaced only the first element: the print gives
% „S e l v kritik“, with „kritik“ set tight. Encoded \emph{Selv}kritik in the
% Danish and \emph{self}-criticism here.
performances. Prof.\ N. has obligingly furnished me with information about my
talent; perhaps a similar piece of information about his own talent will not be
useless to him. Prof.\ N.'s scientific talent is a talent for the special
investigation, where the view of the totality and the fundamental thoughts are
\emph{given} him. Scientific productivity in the higher sense --- the capacity to
bring forth the great, transforming thoughts --- he lacks entirely, and in the
same degree the capacity to master and shape a whole. If he attempts the first,
he attains only to repetitions or bunglings; and if he attempts the last, the
whole at once dissolves for him into an ἄπειρον of unmastered detail. Therefore
Prof.\ N. should never have
% Greek, antiqua italic: the alpha carries the combined smooth breathing and
% acute, ἄ. Copied verbatim.
undertaken a work like the Logic of the Fundamental Ideas, for it \emph{had} to
fail for him, because it demands both the capacities that Prof.\ N. so completely
lacks. How he wants the capacity to delimit and shape a material one gets a
striking example of in the second part of the Logic of the Fundamental Ideas. In
the preface to the first part Prof.\ N. states that it contains ``not quite
two-thirds of the first ninth of the whole work.'' One should then expect that
what was wanting to fill out this ``not quite'' lay tolerably well surveyed
before Prof.\ N.'s gaze. But as he sets about working it up, it swells for him
into so formless a mass that, in order to get that supplement furnished, there
will be needed, besides the whole second part, two more parts of 30 sheets each,
if the same measure is to be kept.

% [text to be added: pp. 23--31 and the Rettelse]

% [text to be added: pp. 23--31 and the Rettelse]

\end{document}
